\documentclass[11pt, titlepage]{article} % DO NOT CHANGE THE FONT SIZE
\usepackage[left=2.2cm,right=2.2cm,top=2.5cm,bottom=2.5cm]{geometry} % DO NOT CHANGE THE MARGINS
\usepackage{amsfonts,amssymb,amsthm,amsmath,amscd}
\usepackage{enumerate}
\usepackage{tikz}
\usepackage{graphicx}
\usepackage{subcaption}
\newcommand{\ecn}{\author}

\title{Appendix: Completely Generated by AI}
\ecn{}
\date{ }


\usepackage{graphicx}

\usepackage[capitalise,noabbrev]{cleveref} % e.g., 
\usepackage{overpic} % use this package for easy labelling of figures
\usepackage{url} % use this package to reference urls, e.g., \url{https://www.maths.cam.ac.uk/postgrad/part-iii/prospective.html}
\usepackage{comment} % use this package to comment things out \begin{comment} blah \end{comment}

\usepackage{algorithm}
\usepackage{algpseudocode}
%FOOTNOTES
\newcommand\blfootnote[1]{% avoids footnotes spilling over
  \begingroup
  \renewcommand\thefootnote{}\footnote{#1}%
  \addtocounter{footnote}{-1}%
  \endgroup
}


\numberwithin{equation}{section}


\newtheorem{theorem}{Theorem}
\newtheorem{lemma}{Lemma}
\newtheorem{corollary}{Corollary}
\newtheorem{proposition}{Proposition}
\newtheorem{definition}{Definition}
\theoremstyle{definition}
\newtheorem{remark}{Remark}
\newtheorem{example}{Example}
\newtheorem*{example*}{Example}
\newtheorem{assumption}{Assumption}

\numberwithin{assumption}{section}
%\numberwithin{theorem}{section}
%\numberwithin{lemma}{section}
%\numberwithin{corollary}{section}
%\numberwithin{proposition}{section}
%\numberwithin{definition}{section}
%\numberwithin{remark}{section}


\begin{document}
\maketitle

\section{Easy proof sketches for the main properties of the MLE}

Let
\begin{equation*}
    \ell_N(\theta)=\log L_N(\theta)=\sum_{i=1}^N \log f(X_i\mid \theta),
\end{equation*}
where $X_1,\dots,X_N \overset{\mathrm{iid}}{\sim} f(x\mid \theta_0)$ and $\theta_0$ is the true parameter.

We assume standard regularity conditions throughout.

\paragraph{Consistency.}

Define the average log-likelihood
\begin{equation*}
    \frac{1}{N}\ell_N(\theta)=\frac{1}{N}\sum_{i=1}^N \log f(X_i\mid \theta).
\end{equation*}
By the law of large numbers, for each fixed $\theta$,
\begin{equation*}
    \frac{1}{N}\ell_N(\theta)\to \mathbb{E}_{\theta_0}[\log f(X\mid \theta)].
\end{equation*}
Now the function
\begin{equation*}
    Q(\theta):=\mathbb{E}_{\theta_0}[\log f(X\mid \theta)]
\end{equation*}
is maximised at $\theta=\theta_0$. Indeed,
\begin{align*}
    Q(\theta_0)-Q(\theta)
    &= \mathbb{E}_{\theta_0}\left[\log \frac{f(X\mid \theta_0)}{f(X\mid \theta)}\right]\\
    &\ge 0,
\end{align*}
with equality only at $\theta=\theta_0$. Therefore the maximiser of the sample log-likelihood should converge to the maximiser of its limit, i.e.
\begin{equation*}
    \hat{\theta}_{\mathrm{MLE}} \to \theta_0.
\end{equation*}
This is the idea behind \textit{consistency}.

\paragraph{Asymptotic normality.}

The MLE satisfies the score equation
\begin{equation*}
    \ell_N'(\hat{\theta}_{\mathrm{MLE}})=0.
\end{equation*}
Now expand $\ell_N'(\hat{\theta}_{\mathrm{MLE}})$ about $\theta_0$:
\begin{equation*}
    0
    =
    \ell_N'(\theta_0)
    +
    \ell_N''(\tilde{\theta}_N)\,
    (\hat{\theta}_{\mathrm{MLE}}-\theta_0),
\end{equation*}
for some $\tilde{\theta}_N$ between $\hat{\theta}_{\mathrm{MLE}}$ and $\theta_0$. Rearranging,
\begin{equation*}
    \hat{\theta}_{\mathrm{MLE}}-\theta_0
    =
    -\frac{\ell_N'(\theta_0)}{\ell_N''(\tilde{\theta}_N)}.
\end{equation*}
Multiply by $\sqrt{N}$:
\begin{equation*}
    \sqrt{N}\,(\hat{\theta}_{\mathrm{MLE}}-\theta_0)
    =
    -\frac{\ell_N'(\theta_0)/\sqrt{N}}{\ell_N''(\tilde{\theta}_N)/N}.
\end{equation*}

Now:
\begin{align*}
    \frac{1}{\sqrt{N}}\ell_N'(\theta_0)
    &= \frac{1}{\sqrt{N}}\sum_{i=1}^N \frac{\partial}{\partial \theta}\log f(X_i\mid \theta_0)
    \overset{d}{\longrightarrow} \mathcal{N}(0,I(\theta_0))
\end{align*}
by the central limit theorem, since the score has mean $0$ and variance $I(\theta_0)$.

Also, by the law of large numbers,
\begin{equation*}
    -\frac{1}{N}\ell_N''(\tilde{\theta}_N)\to I(\theta_0).
\end{equation*}
Hence
\begin{equation*}
    \sqrt{N}\,(\hat{\theta}_{\mathrm{MLE}}-\theta_0)
    \overset{d}{\longrightarrow}
    \mathcal{N}(0,I(\theta_0)^{-1}).
\end{equation*}
Equivalently,
\begin{equation*}
    \hat{\theta}_{\mathrm{MLE}}
    \approx
    \mathcal{N}\!\left(\theta_0,\frac{1}{N}I(\theta_0)^{-1}\right)
\end{equation*}
for large $N$.

\paragraph{Asymptotic unbiasedness.}

From asymptotic normality,
\begin{equation*}
    \hat{\theta}_{\mathrm{MLE}}
    \approx
    \mathcal{N}\!\left(\theta_0,\frac{1}{N}I(\theta_0)^{-1}\right).
\end{equation*}
A normal distribution centred at $\theta_0$ has mean $\theta_0$, so heuristically
\begin{equation*}
    \mathbb{E}[\hat{\theta}_{\mathrm{MLE}}]\to \theta_0.
\end{equation*}
Thus the MLE is \textit{asymptotically unbiased}. 

More precisely, its bias is typically of smaller order than its standard error:
\begin{equation*}
    \mathrm{Bias}(\hat{\theta}_{\mathrm{MLE}})=o(N^{-1/2}).
\end{equation*}

\paragraph{Asymptotic efficiency.}

For any unbiased estimator $T_N$, the Cram\'er--Rao lower bound says
\begin{equation*}
    \mathrm{Var}(T_N)\ge \frac{1}{N\,I(\theta_0)}.
\end{equation*}
But from asymptotic normality,
\begin{equation*}
    \mathrm{Var}(\hat{\theta}_{\mathrm{MLE}})
    \approx
    \frac{1}{N\,I(\theta_0)}.
\end{equation*}
So asymptotically the MLE reaches the lower bound. This is what it means for the MLE to be \textit{asymptotically efficient}.

\paragraph{Functional invariance.}

Let $\alpha=g(\theta)$, and suppose $g$ is a function. The likelihood for $\alpha$ is just the same likelihood written in the transformed parameter:
\begin{equation*}
    L^*(\alpha)=L(\theta)
    \qquad \text{with } \alpha=g(\theta).
\end{equation*}
If $\hat{\theta}_{\mathrm{MLE}}$ maximises $L(\theta)$, then the corresponding transformed value $g(\hat{\theta}_{\mathrm{MLE}})$ maximises the transformed likelihood over $\alpha$. Therefore
\begin{equation*}
    \hat{\alpha}_{\mathrm{MLE}}=g(\hat{\theta}_{\mathrm{MLE}}).
\end{equation*}
So the MLE is \textit{functionally invariant}.

\paragraph{Observed Fisher information.}

Near the MLE, Taylor expansion gives
\begin{equation*}
    \ell_N(\theta)
    \approx
    \ell_N(\hat{\theta}_{\mathrm{MLE}})
    -\frac{1}{2}\hat{I}_N(\theta-\hat{\theta}_{\mathrm{MLE}})^2,
\end{equation*}
where
\begin{equation*}
    \hat{I}_N
    =
    -\ell_N''(\hat{\theta}_{\mathrm{MLE}})
\end{equation*}
is the \textit{observed Fisher information}. Thus the likelihood is approximately Gaussian near its maximum, and
\begin{equation*}
    \mathrm{Var}(\hat{\theta}_{\mathrm{MLE}})
    \approx
    \hat{I}_N^{-1}.
\end{equation*}
Since $\hat{I}_N/N \to I(\theta_0)$, this agrees with the asymptotic variance above.

\section{Example of Bayesian Inference: Gaia Parallax}

For a star at true distance $r>0$, the true parallax is $1/r$ in the usual astronomical units. If the measured parallax is $\varpi$ with known standard deviation $\sigma_\varpi$, we model
\begin{equation*}
    P(\varpi\mid r)
    =
    \frac{1}{\sigma_\varpi\sqrt{2\pi}}
    \exp\!\left[
    -\frac{1}{2\sigma_\varpi^2}
    \left(\varpi-\frac{1}{r}\right)^2
    \right],
    \qquad \sigma_\varpi>0.
\end{equation*}
Thus the likelihood is Gaussian in $\varpi$, but not in $r$. It is convenient to define the \textit{fractional parallax error}
\begin{equation*}
    f=\frac{\sigma_\varpi}{\varpi}.
\end{equation*}

\subsubsection*{Likelihood as a function of distance}

As a function of $r$, the likelihood is
\begin{equation*}
    L(r)=P(\varpi\mid r)
    =
    \frac{1}{\sigma_\varpi\sqrt{2\pi}}
    \exp\!\left[
    -\frac{1}{2\sigma_\varpi^2}
    \left(\varpi-\frac{1}{r}\right)^2
    \right].
\end{equation*}
This is a likelihood in $r$, not a probability density over $r$. The Bayesian posterior is
\begin{equation*}
    P(r\mid \varpi)\propto P(\varpi\mid r)\,P(r),
\end{equation*}
so inference requires a prior on $r$.

Because $\varpi$ is modeled as Gaussian, negative measured parallaxes are possible. These are not meaningless observations; they are noisy measurements which still contain information about $r$.

\subsubsection*{Improper positive prior}

A minimal physical restriction is positivity of distance. Consider the improper prior
\begin{equation*}
    P^*(r)=
    \begin{cases}
        1, & r>0,\\
        0, & \text{otherwise}.
    \end{cases}
\end{equation*}
Then the unnormalized posterior is
\begin{equation*}
    P^*(r\mid \varpi)=
    \begin{cases}
        P(\varpi\mid r), & r>0,\\
        0, & \text{otherwise}.
    \end{cases}
\end{equation*}
This posterior is improper, because
\begin{equation*}
    \lim_{r\to\infty} P^*(r\mid \varpi)=\text{constant}.
\end{equation*}
Hence it is not normalizable, so it has no posterior mean, variance, median, or quantiles.

For $\varpi>0$, the mode is $r=1/\varpi$. For $\varpi\le 0$, there is no finite mode. Thus the likelihood alone does not provide a satisfactory distance estimator.

\subsubsection*{Proper uniform prior with cutoff}

A simple modification is to truncate the prior at some maximum plausible distance $r_{\mathrm{lim}}$:
\begin{equation*}
    P(r)=
    \begin{cases}
        \dfrac{1}{r_{\mathrm{lim}}}, & 0<r\le r_{\mathrm{lim}},\\
        0, & \text{otherwise}.
    \end{cases}
\end{equation*}
Then
\begin{equation*}
    P^*(r\mid \varpi)=
    \begin{cases}
        P(\varpi\mid r), & 0<r\le r_{\mathrm{lim}},\\
        0, & \text{otherwise}.
    \end{cases}
\end{equation*}
This posterior is proper after normalization.

If the posterior mode is used as the estimator, then
\begin{equation*}
    r_{\mathrm{mode}}=
    \begin{cases}
        \dfrac{1}{\varpi}, & 0<\dfrac{1}{\varpi}\le r_{\mathrm{lim}},\\[1ex]
        r_{\mathrm{lim}}, & \dfrac{1}{\varpi}>r_{\mathrm{lim}},\\[0.5ex]
        r_{\mathrm{lim}}, & \varpi\le 0.
    \end{cases}
\end{equation*}
Thus the truncation avoids the divergence to $\infty$, but introduces a hard edge at $r_{\mathrm{lim}}$.

The deeper issue is that a prior uniform in $r$ is not really uninformative. Since $V\propto r^3$, one has $P(V)\,dV=P(r)\,dr$, so $P(r)\propto 1$ implies $P(V)\propto 1/r^2$. Hence a prior uniform in distance corresponds to a stellar volume density proportional to $1/r^2$, which is not astrophysically reasonable.

\subsubsection*{Constant volume-density prior}

A more natural prior is one corresponding to a constant volume density of stars up to $r_{\mathrm{lim}}$:
\begin{equation*}
    P(r)=
    \begin{cases}
        \dfrac{3}{r_{\mathrm{lim}}^3}r^2, & 0<r\le r_{\mathrm{lim}},\\
        0, & \text{otherwise}.
    \end{cases}
\end{equation*}
Then
\begin{equation*}
    P(r\mid \varpi)\propto P(\varpi\mid r)\,P(r),
\end{equation*}
so the unnormalized posterior is
\begin{equation*}
    P_{r^2}^*(r\mid \varpi,\sigma_\varpi)=
    \begin{cases}
        \dfrac{r^2}{\sigma_\varpi}
        \exp\!\left[
        -\dfrac{1}{2\sigma_\varpi^2}
        \left(\varpi-\dfrac{1}{r}\right)^2
        \right], & 0<r\le r_{\mathrm{lim}},\\
        0, & \text{otherwise}.
    \end{cases}
\end{equation*}
This prior is proper and has a clearer physical interpretation.

\subsubsection*{Turning points under the $r^2$ prior}

To locate stationary points of the posterior, differentiate with respect to $r$. This gives two candidate roots,
\begin{equation*}
    r_{\min},\,r_{\mathrm{mode}}
    =
    \frac{1}{\varpi}\frac{1}{4f^2}
    \left(1\pm \sqrt{1-8f^2}\right).
\end{equation*}
The smaller root is the lower mode, and the larger root is a local minimum. These are defined only when
\begin{equation*}
    f<\frac{1}{\sqrt{8}}.
\end{equation*}

This leads to the rule
\begin{equation*}
    r_{\mathrm{est}}=
    \begin{cases}
        r_{\mathrm{mode}}, & \varpi>0,\ f<1/\sqrt{8},\ r_{\mathrm{mode}}\le r_{\mathrm{lim}},\\
        r_{\mathrm{lim}}, & \varpi>0,\ f<1/\sqrt{8},\ r_{\mathrm{mode}}>r_{\mathrm{lim}},\\
        r_{\mathrm{lim}}, & \varpi>0,\ f\ge 1/\sqrt{8},\\
        r_{\mathrm{lim}}, & \varpi\le 0.
    \end{cases}
\end{equation*}
So this prior yields a valid posterior for all $\varpi$, but still inherits the hard cutoff at $r_{\mathrm{lim}}$.

\subsubsection*{Exponentially decreasing volume-density prior}

To remove the hard edge, replace the truncated prior by one with an exponentially decreasing volume density:
\begin{equation*}
    P(r)=
    \begin{cases}
        \dfrac{1}{2L^3}r^2e^{-r/L}, & r>0,\\
        0, & \text{otherwise},
    \end{cases}
\end{equation*}
where $L>0$ is a length scale.

Then
\begin{equation*}
    P(r\mid \varpi)\propto P(\varpi\mid r)\,P(r),
\end{equation*}
and the unnormalized posterior becomes
\begin{equation*}
    P_{r^2e^{-r}}^*(r\mid \varpi,\sigma_\varpi)
    =
    \begin{cases}
        \dfrac{r^2e^{-r/L}}{\sigma_\varpi}
        \exp\!\left[
        -\dfrac{1}{2\sigma_\varpi^2}
        \left(\varpi-\dfrac{1}{r}\right)^2
        \right], & r>0,\\
        0, & \text{otherwise}.
    \end{cases}
\end{equation*}
This posterior is smooth and has no discontinuity at a finite cutoff.

\subsubsection*{Mode equation for the exponential prior}

To find the mode, differentiate the posterior and set the derivative equal to zero. One obtains the cubic equation
\begin{equation*}
    \frac{r^3}{L}-2r^2+\frac{\varpi}{\sigma_\varpi^2}r-\frac{1}{\sigma_\varpi^2}=0.
\end{equation*}
Depending on the values of $\varpi$, $\sigma_\varpi$, and $L$, this cubic may have either one real root or three real roots. Hence the posterior may be either unimodal or bimodal.

A practical rule is as follows:
\begin{enumerate}
    \item If there is one real root, use it as the mode.
    \item If there are three real roots, there are two local maxima:
    \begin{itemize}
        \item if $\varpi\ge 0$, choose the smaller positive mode;
        \item if $\varpi<0$, choose the positive mode.
    \end{itemize}
\end{enumerate}
Thus the exponential prior provides a smooth estimator for all measured parallaxes, including negative ones.

\subsubsection*{Interpretation}

The parallax example shows that the posterior can change substantially with the prior when the data are weak, that is, when $f$ is not small. Priors are therefore not optional decorations: they encode the constraints needed to convert the likelihood into a valid posterior over distance.

The main mathematical lesson is that
\begin{equation*}
    P(r\mid \varpi)\propto P(\varpi\mid r)\,P(r)
\end{equation*}
must be treated as an inference problem in $r$, not as a change-of-variables problem from $\varpi$ to $r$. In particular, simply inverting a noisy parallax is not a valid Bayesian estimator.

\section{Lemma in self-normalised Importance weights}
\begin{lemma}
Let
\begin{equation*}
    \hat I
    =
    \frac{\sum_{i=1}^m f(\theta_i)\tilde w_i}{\sum_{i=1}^m \tilde w_i},
    \qquad
    \tilde w_i=\frac{\tilde P(\theta_i)}{Q(\theta_i)},
\end{equation*}
where $\theta_1,\dots,\theta_m \overset{\mathrm{iid}}{\sim} Q$. Assume the relevant moments exist. Then
\begin{equation*}
    \mathbb{E}[\hat I]=I+\mathcal{O}(m^{-1}),
\end{equation*}
where
\begin{equation*}
    I=\frac{\int f(\theta)\tilde P(\theta)\,d\theta}{\int \tilde P(\theta)\,d\theta}.
\end{equation*}
\end{lemma}

\begin{proof}
Define
\begin{equation*}
    X_i:=f(\theta_i)\tilde w_i,
    \qquad
    Y_i:=\tilde w_i,
\end{equation*}
and the sample means
\begin{equation*}
    \bar X_m:=\frac{1}{m}\sum_{i=1}^m X_i,
    \qquad
    \bar Y_m:=\frac{1}{m}\sum_{i=1}^m Y_i.
\end{equation*}
Then
\begin{equation*}
    \hat I=\frac{\bar X_m}{\bar Y_m}.
\end{equation*}

Let
\begin{equation*}
    \mu_X:=\mathbb{E}_Q[X_i],
    \qquad
    \mu_Y:=\mathbb{E}_Q[Y_i].
\end{equation*}
By definition of $\tilde w_i$,
\begin{align*}
    \mu_X
    &= \int f(\theta)\frac{\tilde P(\theta)}{Q(\theta)}Q(\theta)\,d\theta\\
    &= \int f(\theta)\tilde P(\theta)\,d\theta,
\end{align*}
and similarly
\begin{align*}
    \mu_Y
    &= \int \frac{\tilde P(\theta)}{Q(\theta)}Q(\theta)\,d\theta\\
    &= \int \tilde P(\theta)\,d\theta.
\end{align*}
Hence
\begin{equation*}
    \frac{\mu_X}{\mu_Y}
    =
    \frac{\int f(\theta)\tilde P(\theta)\,d\theta}{\int \tilde P(\theta)\,d\theta}
    = I.
\end{equation*}

Now define
\begin{equation*}
    g(x,y)=\frac{x}{y}.
\end{equation*}
Then $\hat I=g(\bar X_m,\bar Y_m)$. We perform a second-order Taylor expansion of $g$ about $(\mu_X,\mu_Y)$:
\begin{align*}
    g(\bar X_m,\bar Y_m)
    &=
    g(\mu_X,\mu_Y)
    +g_x(\mu_X,\mu_Y)(\bar X_m-\mu_X)
    +g_y(\mu_X,\mu_Y)(\bar Y_m-\mu_Y)\\
    &\quad
    +\frac{1}{2}g_{xx}(\mu_X,\mu_Y)(\bar X_m-\mu_X)^2
    +g_{xy}(\mu_X,\mu_Y)(\bar X_m-\mu_X)(\bar Y_m-\mu_Y)\\
    &\quad
    +\frac{1}{2}g_{yy}(\mu_X,\mu_Y)(\bar Y_m-\mu_Y)^2
    +R_m.
\end{align*}
Since
\begin{equation*}
    g_x(x,y)=\frac{1}{y},
    \qquad
    g_y(x,y)=-\frac{x}{y^2},
    \qquad
    g_{xx}(x,y)=0,
    \qquad
    g_{xy}(x,y)=-\frac{1}{y^2},
    \qquad
    g_{yy}(x,y)=\frac{2x}{y^3},
\end{equation*}
we obtain
\begin{align*}
    \hat I
    &=
    \frac{\mu_X}{\mu_Y}
    +\frac{1}{\mu_Y}(\bar X_m-\mu_X)
    -\frac{\mu_X}{\mu_Y^2}(\bar Y_m-\mu_Y)\\
    &\quad
    -\frac{1}{\mu_Y^2}(\bar X_m-\mu_X)(\bar Y_m-\mu_Y)
    +\frac{\mu_X}{\mu_Y^3}(\bar Y_m-\mu_Y)^2
    +R_m.
\end{align*}

Taking expectations, the linear terms vanish because
\begin{equation*}
    \mathbb{E}[\bar X_m-\mu_X]=0,
    \qquad
    \mathbb{E}[\bar Y_m-\mu_Y]=0.
\end{equation*}
Therefore
\begin{align*}
    \mathbb{E}[\hat I]
    &=
    \frac{\mu_X}{\mu_Y}
    -\frac{1}{\mu_Y^2}\mathbb{E}\!\left[(\bar X_m-\mu_X)(\bar Y_m-\mu_Y)\right]
    +\frac{\mu_X}{\mu_Y^3}\mathbb{E}\!\left[(\bar Y_m-\mu_Y)^2\right]
    +\mathbb{E}[R_m].
\end{align*}

Now
\begin{equation*}
    \mathbb{E}\!\left[(\bar X_m-\mu_X)(\bar Y_m-\mu_Y)\right]
    =\mathrm{Cov}(\bar X_m,\bar Y_m)
    =\frac{1}{m}\mathrm{Cov}(X_1,Y_1),
\end{equation*}
and
\begin{equation*}
    \mathbb{E}\!\left[(\bar Y_m-\mu_Y)^2\right]
    =\mathrm{Var}(\bar Y_m)
    =\frac{1}{m}\mathrm{Var}(Y_1).
\end{equation*}
Hence
\begin{align*}
    \mathbb{E}[\hat I]
    &=
    \frac{\mu_X}{\mu_Y}
    -\frac{1}{m\mu_Y^2}\mathrm{Cov}(X_1,Y_1)
    +\frac{\mu_X}{m\mu_Y^3}\mathrm{Var}(Y_1)
    +\mathbb{E}[R_m].
\end{align*}
Since $\mu_X/\mu_Y=I$, this becomes
\begin{equation*}
    \mathbb{E}[\hat I]
    =
    I
    +\frac{1}{m}\left(
    \frac{\mu_X}{\mu_Y^3}\mathrm{Var}(Y_1)
    -\frac{1}{\mu_Y^2}\mathrm{Cov}(X_1,Y_1)
    \right)
    +\mathbb{E}[R_m].
\end{equation*}
Under the assumed finite-moment conditions, the remainder satisfies
\begin{equation*}
    \mathbb{E}[R_m]=o(m^{-1}).
\end{equation*}
Therefore
\begin{equation*}
    \mathbb{E}[\hat I]=I+\mathcal{O}(m^{-1}),
\end{equation*}
as required.
\end{proof}


\section{Example: Supernova cosmology}
The dimensionless comoving distance to an object with redshift $z$ is
\begin{equation*}
    \tilde d(z;\Omega_M,\Omega_\Lambda,w)
    =
    \int_0^z
    \frac{dz'}
    {\sqrt{\Omega_M(1+z')^3+\Omega_k(1+z')^2+\Omega_\Lambda(1+z')^{3(w+1)}}},
\end{equation*}
where
\begin{equation*}
    \Omega_k=1-\Omega_M-\Omega_\Lambda,
\end{equation*}
and $w$ is the equation-of-state parameter of dark energy.


The theoretical cosmological distance to an object with redshift $z$ inferred from the standard candle method are given by the luminosity distance $d_L$. The dimensionless luminosity distance is
\begin{equation*}
    \tilde d_L(z;\Omega_M,\Omega_\Lambda,w)
    =
    (1+z)
    \begin{cases}
        |\Omega_k|^{-1/2}\sinh\!\bigl(\sqrt{|\Omega_k|}\,\tilde d(z;\Omega)\bigr), & \Omega_k>0,\\
        \tilde d(z;\Omega), & \Omega_k=0,\\
        |\Omega_k|^{-1/2}\sin\!\bigl(\sqrt{|\Omega_k|}\,\tilde d(z;\Omega)\bigr), & \Omega_k<0.
    \end{cases}
\end{equation*}
For a given value of the Hubble constant $H_0$, the dimensionful luminosity distance is
\begin{equation*}
    d_L(z;H_0,\Omega)=\frac{c}{H_0}\tilde d_L(z;\Omega).
\end{equation*}
For our purposes, we use an idealised standard-candle model. Let $M_s$ denote the absolute magnitude of supernova $s$, and let $m_s$ and $z_s$ denote its apparent magnitude and redshift. Assume
\begin{equation*}
    M_s \sim \mathcal{N}(M_0,\sigma_{\mathrm{int}}^2),
\end{equation*}
equivalently,
\begin{equation*}
    M_s=M_0+\varepsilon_{\mathrm{int},s},
    \qquad
    \varepsilon_{\mathrm{int},s}\sim \mathcal{N}(0,\sigma_{\mathrm{int}}^2).
\end{equation*}
Let $\Omega=(\Omega_M,\Omega_L,w)$. The distance modulus defined by
\begin{equation*}
    \mu(z;H_0,\Omega)
    =
    25+5\log_{10}\!\left(\frac{c}{H_0}\tilde d_L(z;\Omega)\right).
\end{equation*}
The apparent magnitude is 
\begin{equation*}
    m_s=M_s+\mu(z_s;H_0,\Omega).
\end{equation*}
Substituting the model for $M_s$ gives
\begin{equation*}
    m_s=M_0+\mu(z_s;H_0,\Omega)+\varepsilon_{\mathrm{int},s}.
\end{equation*}
Hence
\begin{equation*}
    P(m_s\mid z_s;M_0,H_0,\Omega,\sigma_{\mathrm{int}}^2)
    =
    \mathcal{N}\!\left(
    m_s \mid M_0+\mu(z_s;H_0,\Omega),\sigma_{\mathrm{int}}^2
    \right).
\end{equation*}

We can reparametrise by defining the dimensionless constant
\begin{equation*}
    h=\frac{H_0}{100\ \mathrm{km\,s^{-1}\,Mpc^{-1}}}.
\end{equation*}
Then
\begin{align*}
    \mu(z_s;H_0,\Omega)
    &=
    25+5\log_{10}\!\left(\frac{c}{H_0}\tilde d_L(z_s;\Omega)\right)\\
    &=
    25+5\log_{10}\!\left(\frac{c}{100}\tilde d_L(z_s;\Omega)\right)-5\log_{10}h.
\end{align*}
Hence
\begin{equation*}
    M_0+\mu(z_s;H_0,\Omega)
    =
    \underbrace{M_0-5\log_{10}h}_{\bar M_0}
    +
    f(z_s;\Omega),
\end{equation*}
where
\begin{equation*}
    f(z_s;\Omega)
    :=
    25+5\log_{10}\!\left(\frac{c}{100}\tilde d_L(z_s;\Omega)\right).
\end{equation*}
Thus, we have
\begin{equation*}
    P(m_s\mid z_s;\bar M_0,\Omega,\sigma_{\mathrm{int}}^2)
    =
    \mathcal{N}\!\left(
    m_s \mid \bar M_0+f(z_s;\Omega),\sigma_{\mathrm{int}}^2
    \right).
\end{equation*}
So $M_0$ and $\log_{10}h$ cannot be separately constrained (not separately identifiable given the data).

\subsubsection*{Bayesian Inference}

Assume the data are independent, the likelihood is
\begin{equation*}
    L(\bar M_0,\Omega,\sigma_{\mathrm{int}}^2)
    =
    \prod_{s=1}^N
    P(m_s\mid z_s;\bar M_0,\Omega,\sigma_{\mathrm{int}}^2).
\end{equation*}
A simple improper prior choice is
\begin{equation*}
    \bar M_0,\Omega_\Lambda \sim U(-\infty,\infty),
    \qquad
    \Omega_M,\sigma_{\mathrm{int}}^2 \sim U(0,\infty),
\end{equation*}
then the posterior is
\begin{equation*}
    P(\bar M_0,\Omega,\sigma_{\mathrm{int}}^2\mid \{m_s,z_s\})
    \propto
    L(\bar M_0,\Omega,\sigma_{\mathrm{int}}^2)\,
    P(\bar M_0)\,P(\Omega_M)\,P(\Omega_\Lambda)\,P(\sigma_{\mathrm{int}}^2).
\end{equation*}

\subsubsection*{Derived quantities}

For the case $w=-1$, an important derived quantity is the deceleration parameter
\begin{equation*}
    q_0=\frac{\Omega_M}{2}-\Omega_\Lambda.
\end{equation*}
Posterior draws of $(\Omega_M,\Omega_\Lambda)$ immediately induce posterior draws of $q_0$. If most of the posterior mass lies below zero, then
\begin{equation*}
    q_0<0,
\end{equation*}
which means the cosmic expansion is accelerating.

\subsubsection*{Flat-universe model with free {$w$}{w}}

A second case study assumes a flat universe,
\begin{equation*}
    \Omega_\Lambda=1-\Omega_M,
\end{equation*}
but now treats $w$ as unknown. The parameter vector becomes
\begin{equation*}
    \theta=(\bar M_0,\sigma_{\mathrm{int}}^2,\Omega_M,w).
\end{equation*}
The same supernova likelihood is used, but now with
$f(z_s;\Omega_M,w)$ computed under the flat-universe restriction.

This gives a four-dimensional posterior, which serves as a useful example for comparing MCMC algorithms.

\subsubsection*{Single-chain trace plots}

A trace plot displays the sampled parameter values against iteration number. In the supernova example, separate trace plots are shown for the four parameters
$\bar M_0$, $\sigma_{\mathrm{int}}$, $\Omega_M$, and $\Omega_\Lambda$ or $w$.
Trace plots help assess:
\begin{itemize}
    \item whether the chain has moved away from its starting point;
    \item whether it seems to fluctuate around a stable region;
    \item whether mixing appears slow;
    \item whether certain parameters are more difficult to sample than others.
\end{itemize}

\subsubsection*{Multiple independent chains}

A stronger diagnostic is to run several overdispersed chains independently. If all chains mix well and appear to explore the same region of the posterior, this provides stronger evidence of convergence than a single trace plot.

The lecture slides illustrate this by showing four independent chains for the cosmology parameters. After burn-in, one may combine all retained samples from all chains and treat them as one posterior sample.

\section{Gibbs}

\paragraph{d-dimensional Gibbs sampler.}
More generally, let $\vec{\theta}=(\theta_1,\dots,\theta_d)$. A Gibbs sampler updates the parameters one at a time within each cycle. At iteration cycle $t$, define the intermediate state before updating the $j$th coordinate by
\begin{equation*}
\vec{\theta}^{(t,j-1)}
:=
\bigl(
\theta_1^{(t)},\dots,\theta_{j-1}^{(t)},\theta_j^{(t-1)},\theta_{j+1}^{(t-1)},\dots,\theta_d^{(t-1)}
\bigr),
\end{equation*}
for $j=1,\dots,d$, with $\vec{\theta}^{(t,0)}:=\vec{\theta}^{(t-1)}$.

We also define
\begin{equation*}
\vec{\theta}_{-j}^{(t,j-1)}
:=
\bigl(
\theta_1^{(t)},\dots,\theta_{j-1}^{(t)},\theta_{j+1}^{(t-1)},\dots,\theta_d^{(t-1)}
\bigr),
\end{equation*}
that is, the current components of $\vec{\theta}$ except $\theta_j$, using the already updated values up to coordinate $j-1$.

We use the conditional posterior distribution for the Gibbs update for the $j$th component:
\begin{equation*}
\theta_j^{(t)} \sim P\bigl(\theta_j \mid \vec{\theta}_{-j}^{(t,j-1)}, D\bigr).
\end{equation*}

Equivalently, the proposal distribution at the $j$th substep only updates the $j$th component of $\vec{\theta}^{(t,j-1)}$, so we can write the proposal distribution as
\begin{equation*}
J\bigl(\vec{\theta}^{\,*} \mid \vec{\theta}^{(t,j-1)}\bigr)
=
P\bigl(\theta_j^* \mid \vec{\theta}_{-j}^{(t,j-1)}, D\bigr)
\,\mathbf{1}\bigl\{\vec{\theta}_{-j}^{\,*}=\vec{\theta}_{-j}^{(t,j-1)}\bigr\}.
\end{equation*}

Note that the indicator makes sure that $J$ only updates the $j$th component of $\vec{\theta}^{(t,j-1)}$. Now, by construction of $J$, we have
\begin{align*}
J\bigl(\vec{\theta}^{(t,j-1)} \mid \vec{\theta}^{\,*}\bigr)
&=
P\bigl(\theta_j^{(t-1)} \mid \vec{\theta}_{-j}^{\,*}, D\bigr)
\,\mathbf{1}\bigl\{\vec{\theta}_{-j}^{(t,j-1)}=\vec{\theta}_{-j}^{\,*}\bigr\} \\
&=
P\bigl(\theta_j^{(t-1)} \mid \vec{\theta}_{-j}^{(t,j-1)}, D\bigr)
\,\mathbf{1}\bigl\{\vec{\theta}_{-j}^{\,*}=\vec{\theta}_{-j}^{(t,j-1)}\bigr\}.
\end{align*}

Substituting this proposal into the Metropolis--Hastings ratio gives
\begin{align*}
r
&=
\frac{
P(\vec{\theta}^{\,*} \mid D)\,
J\bigl(\vec{\theta}^{(t,j-1)} \mid \vec{\theta}^{\,*}\bigr)
}{
P\bigl(\vec{\theta}^{(t,j-1)} \mid D\bigr)\,
J\bigl(\vec{\theta}^{\,*} \mid \vec{\theta}^{(t,j-1)}\bigr)
} \\
&=
\frac{
P\bigl(\theta_j^*,\vec{\theta}_{-j}^{(t,j-1)} \mid D\bigr)\,
P\bigl(\theta_j^{(t-1)} \mid \vec{\theta}_{-j}^{(t,j-1)}, D\bigr)
}{
P\bigl(\theta_j^{(t-1)},\vec{\theta}_{-j}^{(t,j-1)} \mid D\bigr)\,
P\bigl(\theta_j^* \mid \vec{\theta}_{-j}^{(t,j-1)}, D\bigr)
} \\
&=
\frac{
P\bigl(\theta_j^* \mid \vec{\theta}_{-j}^{(t,j-1)}, D\bigr)\,
P\bigl(\vec{\theta}_{-j}^{(t,j-1)} \mid D\bigr)\,
P\bigl(\theta_j^{(t-1)} \mid \vec{\theta}_{-j}^{(t,j-1)}, D\bigr)
}{
P\bigl(\theta_j^{(t-1)} \mid \vec{\theta}_{-j}^{(t,j-1)}, D\bigr)\,
P\bigl(\vec{\theta}_{-j}^{(t,j-1)} \mid D\bigr)\,
P\bigl(\theta_j^* \mid \vec{\theta}_{-j}^{(t,j-1)}, D\bigr)
} \\
&= 1.
\end{align*}

Hence each Gibbs update is always accepted. After updating all coordinates $j=1,\dots,d$, we set
\begin{equation*}
\vec{\theta}^{(t)}:=\vec{\theta}^{(t,d)}.
\end{equation*}
\end{document}